\def\topic{}
\ifdefined\MyHeaderIncluded
\else
\def\qwertyflag{true}

\ifdefined\MyHeaderIncluded

\else
  \newcommand{\foobar}{FooBar}
  \def\MyHeaderIncluded{true}
% if is till the end of the file

\documentclass[10pt,a4paper]{article}
\usepackage[utf8]{inputenc}
%\usepackage[german]{babel}
%\usepackage[cyrillic]{babel}
\usepackage[T1]{fontenc}
%\usepackage[T2A]{fontenc}
%\usepackage[english, russian]{babel}
\usepackage{amsmath}
\usepackage{amsfonts}
\usepackage{amssymb}
\usepackage{amsthm}
\usepackage{mathtools} %coloneqq

\usepackage{bussproofs} %

%\renewcommand{\thesubsection}{\arabic{subsection}}
%\newtheorem{thm}{Theorem}[section]
%\renewcommand{\thethm}{\arabic{subsection}.\arabic{thm}}

\usepackage{thmtools}
\usepackage{hyperref}
\hypersetup{
    colorlinks=true,
    linkcolor=blue,
    filecolor=magenta,      
    urlcolor=cyan,
    pdftitle={Overleaf Example},
    %pdfpagemode=FullScreen,
    }
%\usepackage{stix-type1}%\turnediota
\theoremstyle{definition}

\newtheorem{constructor}{Constructor}[section]
\newtheorem{iaxiom}{Axiom}
\newtheorem{axiom}{Axiom}
\newtheorem*{axiom*}{Axiom}
\newtheorem{definition}{Definition}[section]
\newtheorem*{definition*}{Definition}
\newtheorem{remark}{Remark}[section]
\newtheorem*{remark*}{Remark}
\newtheorem{recall}{Recall}[section]
\newtheorem*{recall*}{Recall}

\newtheorem{lemma}{Lemma}[section]
\newtheorem{theorem}{Theorem}[section]
\newtheorem*{theorem*}{Theorem}
\newtheorem{note}{Note}[section]
\usepackage{etoolbox}

%\newcommand*{\@Dashv}[2]{%
%  \reflectbox{$\m@th#1#2$}%
%}
\newcommand*{\separation}[3]{% 1)binder, 2)set, 3)condition
  \{#1\in #2: #3\}
}

% https://www.freecodecamp.org/news/simple-regex-tricks-for-beginners-3acb3fa257cb/
%  \\mathcal\{T\} |---> \topology
%  \\mathcal\{([^\}]*)\}   |--->  $1
\newcommand{\AND}{\ \ \&\ \ }
\newcommand{\any}{\mathrm{\ any}}
\newcommand{\borel}{\mathcal{B}}
\newcommand{\can}{\mathrm{can}}
\newcommand{\canonize}{\mathrm{canonize}}
\newcommand{\closure}{\mbox{cl}}
\newcommand{\complex}{\mathbb{C}}
\newcommand{\converges}{\longrightarrow}
\newcommand{\countable}{\mbox{\ countable}}
\newcommand{\ldefined}{\mbox{\ defined}}
\newcommand{\dom}{\mbox{dom}}
\newcommand{\Dom}{\mbox{Dom}}
\newcommand{\Done}{\noindent \textbf{Done.}\\}
\newcommand{\lexists}{\mbox{\ exists}}
\newcommand{\lforall}{\mbox{\ forall}}
\newcommand{\lunique}{\mbox{\ unique}}
\newcommand{\family}{\mathrm{\ family}}
\newcommand{\Fm}{\mathrm{Fm}}
\newcommand{\FuncSymb}{\mathrm{FuncSymb}}
\newcommand{\function}{\mathrm{\ function}}
\newcommand{\functional}{\mbox{\ functional}}
\newcommand{\IFF}{\ \Longleftrightarrow\ }
\renewcommand{\iff}{\ \longleftrightarrow\ }
\newcommand{\IMPLIES}{\ \Longrightarrow\ }
\renewcommand{\implies}{\ \longrightarrow\ }
\newcommand{\Inf}{\mathrm{Inf}}
\newcommand{\injective}{\mbox{\ injective}}
\newcommand{\injection}{\mbox{\ injection}}
\newcommand{\interior}{\mbox{int}}
\newcommand{\intersection}{\bigcap}
\newcommand{\isempty}{\mathrm{\ empty}}
\newcommand{\LB}{\mathrm{LB}}
\newcommand{\nghbset}{\mathcal{N}}
\newcommand{\nonempty}{\mathrm{\ nonempty}}
\newcommand{\Ord}{\mathrm{Ord}}
\newcommand{\power}{\mathcal{P}}
\newcommand{\pr}{\mathrm{pr}}
\newcommand{\fst}{\mathrm{fst}}
\newcommand{\snd}{\mathrm{snd}}
\newcommand{\PredSymb}{\mathrm{PredSymb}}
\newcommand{\Proof}{\noindent \textbf{Proof.}\\}
\newcommand{\ran}{\mathrm{ran}}
\newcommand{\Ran}{\mathrm{Ran}}
\newcommand{\real}{\mathbb{R}}
\newcommand{\relation}{\mathrm{\ relation}}
\newcommand{\set}{\mathrm{\ set}}
\newcommand{\singleton}{\mathrm{\ singleton}}
\newcommand{\semantics}{\mathrm{semantics}}
\newcommand{\semialgebra}{\mathcal{S}}
\newcommand{\Sup}{\mathrm{Sup}}
\newcommand{\surjective}{\mbox{\ surjective\ on\ }}
\newcommand{\Term}{\mathrm{Term}}
\newcommand{\Type}{\mathrm{Type}}
\newcommand{\topology}{\mathcal{T}}
\newcommand{\total}{\mbox{\ total\ on\ }}
\newcommand{\transitive}{\mathrm{\ transitive}}
\newcommand{\union}{\bigcup}
\newcommand{\lTrue}{\mathrm{True}}
\newcommand{\lFalse}{\mathrm{False}}
\newcommand{\lUncertain}{\mathrm{Uncertain}}
\newcommand{\lif}{\mathrm{if}}
\newcommand{\lthen}{\mathrm{then}}
\newcommand{\lelse}{\mathrm{else}}
\newcommand{\lexception}{\mathrm{exception}}
\newcommand{\lundef}{\mathrm{undef}}
\newcommand{\univocal}{\mathrm{\ univocal}}
\newcommand{\univocalat}{\mathrm{\ univocal\ at\ }}
\newcommand{\UB}{\mathrm{UB}}
\newcommand{\Un}{\mathcal{U}}
\newcommand{\Var}{\mathrm{Var}}

\newcommand{\notvarphi}{\phi}
\renewcommand{\phi}{\varphi}

\newcommand{\comment}[1]{\quad\mbox{(#1)}}
% local commands
%\preto\section{\undef\blipo} %etoolbox
\preto\subsection{\undef\blipo} %etoolbox
% https://tex.stackexchange.com/questions/105684/limit-macro-newcommand-scope-to-section
\usepackage{graphicx}
\usepackage{authblk}

\setlength\parindent{0pt} % whole doc of \noindent
\usepackage{soul}
%\usepackage{mathgl}

\fi % guard
\ifdefined\topic
\else
\def\topic{No Titel}
\fi
\title{\topic}
\author[1]{Georgy Dunaev}
\begin{document}
\maketitle
\fi
\section{Topology}
%(\cite[p.450]{Simon1982})
%(by \cite[p.469]{Treves1967})
This chapter explores several descriptions of topology, inspired by
\emph{Topological Vector Spaces, Distributions and Kernels}
(\cite{Treves1967}). The point--neighbourhood relation, families of basic
sets, and families of open sets provide different starting points.
The conventions are stated locally: in particular, the first use of
``subbase'' means an intersection-closed family, while the later
``generator'' need not be intersection-closed.
\subsection{Filter}
\begin{definition}[filter]
%\[\]
The set $F$ is called a \textit{filter} in the local convention below if
it satisfies the following conditions. Its carrier is taken to be
$\union F$, rather than supplied separately.
\[
\left( \emptyclass\notin F \vphantom\union \right),
\tag{F$_1$}
\]
\[
\left(
\middle(
A\in F
\middle)
\land
\middle(
B\in F
\middle)
\middle)
\implies
\middle(\middle(A\cap B\middle)\in F \vphantom\union \right),
\tag{F$_2$}
\]
\[
\left(F \ni A \subseteq B \subseteq \union F\middle)
\implies
\middle(B\in F\right).
\tag{F$_3$}
\]

\iffalse
\begin{align*}
\left( \emptyclass\notin F \vphantom\union \right),
 && \text{(F$_1$)}\\
\left(
\middle(
A\in F
\middle)
\land
\middle(
B\in F
\middle)
\middle)
\implies
\middle(\middle(A\cap B\middle)\in F \vphantom\union \right),
 && \text{(F$_2$)}\\
\left(F \ni A \subseteq B \subseteq \union F\middle)
\implies
\middle(B\in F\right).
&& \text{(F$_3$)}\\
\end{align*}
\fi
%    x &=4-3 && \text{Subtract 3 from both sides}\\
%    x &=1   && x \text{ must be one}
\end{definition}
These three conditions also admit $F=\emptyclass$, because the last two
are then vacuous. Thus this convention includes an empty family. When a
nonempty filter is required, that hypothesis will be stated separately;
for a nonempty $F$, upward closure implies $\union F\in F$.
\subsection{Filter Base}
\begin{definition}[filter base with its recovered carrier]
A filter base is a nonempty set $\mathcal B$ of nonempty sets such that
for $B_1,B_2\in\mathcal B$ there is $B_3\in\mathcal B$ with
$B_3\subseteq B_1\cap B_2$. Its carrier is $X=\bigcup\mathcal B$.
Define the generated family by
\[
F(\mathcal B)=\{A\subseteq X:\exists B\in\mathcal B,\ B\subseteq A\}.
\]
Thus the base need not itself be upward closed or closed under
intersections. The refinement condition supplies a member inside each
intersection of two members.
\end{definition}

\begin{proposition}[Generation and recovery of a nonempty filter]
\label{prop:filter_base_generation}
The family $F(\mathcal B)$ is a nonempty filter in the preceding local
convention, has carrier $X$, and is the least such filter containing
$\mathcal B$. Conversely, every nonempty filter $F$ in that convention
is a filter base and satisfies $F(F)=F$.
\end{proposition}
\Proof
Every base member belongs to $F(\mathcal B)$, so the generated family
is nonempty and its union is $X$: containment in $X$ follows from its
definition, and the reverse containment follows from
$\mathcal B\subseteq F(\mathcal B)$. No member is empty, since it
contains a nonempty base member. If $A_1,A_2\in F(\mathcal B)$,
choose $B_i\in\mathcal B$ with $B_i\subseteq A_i$ for $i=1,2$.
A common refinement $B_3$ is contained in $A_1\cap A_2$, so this
intersection belongs to $F(\mathcal B)$. Upward closure within $X$
is immediate. These are exactly the three filter axioms.

If a filter $G$ contains $\mathcal B$, then $X\subseteq\bigcup G$.
For $B\subseteq A\subseteq X$ with $B\in\mathcal B$, upward
closure of $G$ therefore gives $A\in G$. Hence
$F(\mathcal B)\subseteq G$, even if $G$ has a larger carrier.
Finally, a nonempty filter is a base by its intersection axiom, and
its upward closure gives $F(F)=F$. The empty filter admitted earlier
is deliberately excluded from this base representation.
\Done

\subsection{Topology via a point--neighbourhood relation}
The relation $x\mathcal{R}S$ is intended to mean that $S$ is a
neighbourhood of $x$. In this set-sized development, take $\mathcal{R}$
to be a set relation so that its images and the families defined below
are sets. The carrier $E$ is recovered from the relation. Read the
point axioms for $x\in E$ and the set axioms for $S\subseteq E$.
The displayed axioms specify this proposed presentation; equivalence
with other presentations requires a separate argument.
\begin{definition}[relation between points and neighbourhoods](parameter)
\[
\mathcal{R}\subseteq \V\times \V
\]
\end{definition}
\begin{definition}[relation between neighbourhoods and points]
\[
\reflectbox{$\mathcal{R}$}\ceq \mathcal{R}^{T}
\]
\end{definition}
\begin{definition}[neighbourhoods of a point]
\[
\mathrm{tni}(x)\ceq \mathcal{R}[\{x\}]
\]
%\[
%\mathcal{N}(x)\ceq \mathcal{R}[\{x\}]
%\]
\end{definition}
\begin{definition}[interior of a set]
\[
\interior(S)\ceq \reflectbox{$\mathcal{R}$}[\{S\}]
\]
%\[
%\reflectbox{$\mathcal{N}$}
%(S)\ceq \reflectbox{$\mathcal{R}$}[\{S\}]
%\]
\end{definition}
\begin{definition}[Candidate family of open sets]
\[
\mathcal{T}\ceq \ran(\interior)
\]
%\[
%\mathcal{T}\ceq \ran(\reflectbox{$\mathcal{N}$})
%\]
\end{definition}
\begin{definition}[Family of neighbourhood filters]
\[
\mathcal{N}\ceq \ran(\mathrm{tni})
\]
%\[
%\mathbf{N}\ceq \ran(\mathcal{N})
%\]
\end{definition}
\begin{definition}[set of all points]
%\mathbf{NN}\ceq \union\ran(\mathcal{N})
\[
E\ceq \union\mathcal{T}
\]
\end{definition}
\begin{definition}[set of all neighbourhoods]
%\mathbf{NN}\ceq \union\ran(\mathcal{N})
\[
N\ceq \union\mathcal{N}
\]
%\[
%\mathbf{NN}\ceq \union\mathbf{N}
%\]
\end{definition}

\begin{axiom*}[N$_0$]
\[
(\mathrm{tni}(x)\ \mathrm{filter})
\]
\end{axiom*}


\begin{axiom*}[N$_1$]
% x\in\intersection \mathcal{N}(x)
\[
x\in\intersection \mathrm{tni}(x)
\]
\end{axiom*}

\begin{axiom*}[N$_2$]
\[
\interior(S)\in \left(N\cup\{\emptyclass\}\right)
\]
\end{axiom*}
\paragraph{Scope of the proposed neighbourhood axioms.}
Condition (N$_2$) says that a nonempty interior is a neighbourhood of
\emph{some} point. This is weaker than requiring it to be a neighbourhood
of every point it contains. The difference affects the intended passage
from the relation to an ordinary topological interior.

For example, let $E=\{a,b,c,d\}$, with four distinct points, and put
\[
U_a=\{a,b,d\},\quad U_b=\{b,c,d\},\quad
U_c=\{c,d\},\quad U_d=\{d\}.
\]
Declare $x\mathcal{R}S$ exactly when $U_x\subseteq S\subseteq E$.
Each neighbourhood family is a nonempty principal filter, and every
neighbourhood of $x$ contains $x$, so (N$_0$) and (N$_1$) hold.
Whenever $\interior(S)$ is nonempty, $d\in S$, hence
$d\in\interior(S)$. Consequently $\interior(S)$ is a neighbourhood
of $d$, proving (N$_2$). Nevertheless,
\[
\interior(\{a,b,d\})=\{a,d\},\qquad
\interior(\{a,d\})=\{d\}.
\]
Thus this interior is not idempotent. The three axioms retain a coherent
weaker setting, but alone do not identify $\interior$ with the interior
operator of an ordinary topology.

For that identification, a natural strengthening is the pointwise clause
\[
x\in\interior(S)\implies
\interior(S)\in\mathrm{tni}(x).
\tag{N$_2^{\mathrm{pt}}$}
\]
Read this with nonempty neighbourhood filters on the common carrier $E$,
and with every neighbourhood a subset of $E$. It says that each
neighbourhood contains an open neighbourhood of the same point.
The subbase construction below satisfies this stronger clause too: its
basic witness $B$ lies inside $\interior(S)$ and contains the original
point $x$. The weaker proposal is therefore preserved without conflating
its scope with that of the ordinary topological setting.

\begin{theorem}[Recovering a topology from the pointwise axioms]
Suppose every $\mathrm{tni}(x)$, $x\in E$, is a nonempty filter on
$E$, every member contains $x$, and (N$_2^{\mathrm{pt}}$) holds.
Put $I(S)=\{x\in E\mid S\in\mathrm{tni}(x)\}$ for $S\subseteq E$.
Then
\[
\mathcal{T}=\{U\subseteq E\mid I(U)=U\}
\]
is a topology on $E$, and its neighbourhood relation is the given one.
\end{theorem}
\Proof
The filter axioms give $I(E)=E$, $I(\emptyclass)=\emptyclass$, and
$I(S\cap T)=I(S)\cap I(T)$. They also give monotonicity:
$S\subseteq T$ implies $I(S)\subseteq I(T)$. The point condition
gives $I(S)\subseteq S$; (N$_2^{\mathrm{pt}}$) gives
$I(S)\subseteq I(I(S))$. Applying the first inclusion to $I(S)$
proves equality, so $I$ is idempotent. Its range is therefore exactly
its family of fixed points, agreeing with the earlier notation
$\mathcal{T}=\ran(I)$.

Its fixed points contain $\emptyclass,E$ and are closed under binary
intersections. If each $U_j$ is fixed, monotonicity gives
$U_j\subseteq I(\bigcup_j U_j)$, so their union is fixed too.
Thus the fixed points form a topology. Finally, if
$S\in\mathrm{tni}(x)$, then $x\in I(S)\subseteq S$ and $I(S)$ is
open. Conversely, an open $U$ with $x\in U\subseteq S$ gives
$U\in\mathrm{tni}(x)$ and hence $S\in\mathrm{tni}(x)$ by upward
closure. This proves equality of the neighbourhood relations.
\Done

\subsection{Topology via (sub)base}
\begin{definition}[subbase]
In this subsection, a \textit{subbase} is a set $\mathcal{B}$ closed
under binary intersections. This is a local naming convention: the
closure condition already provides the refinement property needed for
a base on the carrier $\union\mathcal{B}$. Empty members are allowed.
\begin{align*}
\left(A\in \mathcal{B}\middle)\land \middle(B\in \mathcal{B}\middle)
\implies
\middle(A\cap B\in \mathcal{B} \vphantom\union \right),
\end{align*}
% && \text{(F2)}\\
\end{definition}
%point--neighbourhood relation defined via topology base?
\begin{definition}[topology via subbase]
%For $p\in E$, $S\subseteq E$, and a set $\mathcal{B}\subseteq\power E$(later we will see that its elements are open), such that $\emptyclass\notin \mathcal{B}$,
For any $p$, any $S$, and a subbase $\mathcal{B}$, we define:
\footnote{
Here we use $\union\mathcal{B}$
%in order to follow the assumption that something generated by a basis is uniquely determined. E.g. ``The filter
%generated by the [filter] basis is uniquely determined'', p.7 TVS.
%Otherwise
, because otherwise
any generated topology would always have a class of all sets as a carrier.
%Otherwise we couldn't distinguish topology from some of its subtopologies.
Later we will also see that elements of $\mathcal{B}$ are open.
}
\[
\left(p \mathcal{R}_\mathcal{B} S \right) :\LRA
\left(\exists B\in\mathcal{B},\  p\in B \subseteq S \subseteq \union\mathcal{B}  \right)
\]
\end{definition}
%(Note:
%)

%point--neighbourhood relation defined via filter bases is a point--neighbourhood relation
\begin{theorem}[topology via subbase is a topology]
The relation $\mathcal{R}_\mathcal{B}$ between points and neighbourhoods is topological.
\ProofNoPar
Put $E=\union\mathcal{B}$ and fix $x\in E$.
\begin{itemize}
\item (N$_0$) Every neighbourhood contains $x$, so the empty set is not
a neighbourhood. If $S$ and $T$ are neighbourhoods, choose
$B_1,B_2\in\mathcal{B}$ with $x\in B_1\subseteq S$ and
$x\in B_2\subseteq T$. Then $B_1\cap B_2\in\mathcal{B}$ and
$x\in B_1\cap B_2\subseteq S\cap T$, so $S\cap T$ is a neighbourhood.
Upward closure within $E$ follows immediately from the definition.
There is some $B\in\mathcal{B}$ containing $x$, so $E$ itself is a
neighbourhood and this neighbourhood family is nonempty.
\item (N$_1$) Each member $S$ of the neighbourhood family contains a
basic set containing $x$, hence contains $x$. Therefore $x$ belongs
to the intersection of the family.
\item (N$_2$) Let $I=\interior_\mathcal{B}(S)$. If $I=\emptyclass$,
the axiom holds. Otherwise take $x\in I$ and choose
$B\in\mathcal{B}$ with $x\in B\subseteq S\subseteq E$.
For every $y\in B$, the same $B$ witnesses $y\mathcal{R}_\mathcal{B}S$,
so $B\subseteq I\subseteq E$. Thus $I$ is a neighbourhood of $x$,
and consequently $I\in N_\mathcal{B}$.
\end{itemize}
\DoneNoPar
\end{theorem}
\begin{definition}[Rectangles from neighbourhoods]
The following family collects neighbourhood rectangles. Its local
subfamilies at a point provide the product filter bases below. Calling
it an open base additionally requires the rectangles to be open.
\[
\mathcal{B}_{E\times F} \ceq \cc{A\times B}{(A\in N_E) \land (B\in N_F)}
\]
\end{definition}
To describe this construction at a specified point, use the notation
\begin{notation}[neighbourhoods of a point]
\[
\mathrm{nghb} \ceq \mathrm{tni}
\]
\end{notation}
\begin{definition}[product filter basis function]
\[
\mathrm{pfb}((a,b))\ceq \cc{A\times B }{ (A\in \mathrm{nghb}_E(a)) \land (B\in \mathrm{nghb}_F(b))}
\]
\end{definition}
%\pagebreak
%\subsection{Filter and filter base}
\pagebreak
\subsection{Topologies and intersection-closed structures}
Here a topology or sigma-algebra carries its own underlying set: the
carrier is the union of its members. This explains why a carrier does
not appear as a separate argument in the definitions.
\begin{definition}[topology]
The set $\mathcal{T}$ is called \textit{topology} iff
%\left( \emptyclass\notin F \vphantom\union \right),
% && \text{(F$_1$)}\\
\begin{align*}
\left(S \subseteq \mathcal{T}\middle)
\implies
\middle(\middle(\union S\middle)\in \mathcal{T} \vphantom\union \right),
 && \text{(T$_1$)}\\
\left(\middle(
A\in \mathcal{T}\middle)\land \middle(B\in \mathcal{T}\middle)\middle)
\implies
\middle(\middle(A\cap B\middle)\in \mathcal{T} \vphantom\union \right),
 && \text{(T$_2$)}\\
\end{align*}
%\left(F \ni A \subseteq B \subseteq \union F\middle)
%\implies
%\middle(B\in F\right).
%&& \text{(F$_3$)}\\
\end{definition}

\begin{definition}[sigma--algebra]
The set $\sigma$ is called \textit{sigma--algebra} iff
\begin{align*}
\left(
\middle(
S \subseteq \sigma
\middle)
\land
\middle(
|S|\leqslant \omega
\middle)
\middle)
\implies
\middle(\middle(\union S\middle)\in \sigma \vphantom\union \right),
 && \text{(S$_1$)}\\
 \left(\middle(
A\in \sigma \middle)\land \middle(B\in \sigma\middle)\middle)
\implies
\middle(\middle(A\setminus B\middle)\in \sigma \vphantom\union \right),
 && \text{(S$_2$)}\\
\left( \left(\union \sigma\right) \in \sigma \vphantom\union \right).
 && \text{(S$_3$)}\\
\end{align*}
%\left(\middle(
%A\in \mathcal{T}\middle)\land \middle(B\in \mathcal{T}\middle)\middle)
%\implies
%\middle(\middle(A\cap B\middle)\in \mathcal{T} \vphantom\union \right),
% && \text{(T$_2$)}\\
\end{definition}
For all the intersection statements below, $Q$ is a \emph{nonempty}
family of the indicated sets. This ensures that $\intersection Q$ is a
set. When the structures are subspaces, closed sets, convex sets, or
balanced sets, they lie in one fixed ambient vector space or topological
space. Empty intersections can instead be treated relative to a stated
ambient set; the unrestricted class convention
$\intersection\emptyclass=\V$ is not used for these statements.

\begin{theorem}[intersection of topologies]
\begin{align*}
\left(
Q \subseteq \cc{\mathcal{T}}{\mathcal{T}\mathrm{\ topology}}
\middle)
\implies
\middle(
\middle(
\intersection Q
\middle)
\mathrm{\ topology}
\vphantom\union
\right)
\end{align*}
% && \text{(S$_2$)}\\
\Proof
Let $\mathcal{I}=\intersection Q$. If $S\subseteq\mathcal{I}$,
then $S\subseteq\mathcal{T}$ for every $\mathcal{T}\in Q$, so
$\union S\in\mathcal{T}$ for every such $\mathcal{T}$. Thus
$\union S\in\mathcal{I}$. Binary intersections are inherited in the
same way. Nonemptiness of $Q$ bounds $\mathcal{I}$ by a set topology,
so $\mathcal{I}$ is a set and satisfies both defining axioms.
\DoneNoPar
\end{theorem}


\begin{theorem}[intersection of sigma--algebras]
\begin{align*}
\left(
\middle(
Q \subseteq \cc{\sigma}{\sigma\mathrm{\ sigma{\shortminus}algebra}}
\middle)
\land
\middle(
\middle(
\union\intersection Q
\middle)
\in \intersection Q
\middle)
\middle)
\implies
\middle(
\middle(
\intersection Q
\middle)
\mathrm{\ sigma{\shortminus}algebra}
\vphantom\union
\right)
\end{align*}
\Proof
Put $\mathcal{I}=\intersection Q$. Countable unions and relative
differences of members of $\mathcal{I}$ belong to every sigma-algebra
in $Q$, hence to $\mathcal{I}$. The additional displayed hypothesis
is precisely the remaining carrier axiom
$\union\mathcal{I}\in\mathcal{I}$.
\DoneNoPar
\end{theorem}


\begin{theorem}[intersection of linear subspaces]
\begin{align*}
\left(
Q \subseteq \cc{L}{L\mathrm{\ linear\ subspace}}
\middle)
\implies
\middle(
\middle(
\intersection Q
\middle)
\mathrm{\ linear\ subspace}
\vphantom\union
\right)
\end{align*}
\Proof
Every member of $Q$ contains the zero vector, so their intersection
does as well. If $v_1,v_2\in\intersection Q$ and
$\alpha_1,\alpha_2\in\mathbb{K}$, then
$\alpha_1v_1+\alpha_2v_2$ lies in every member of $Q$, hence in their
intersection. Thus the intersection is a linear subspace.
\DoneNoPar
\end{theorem}

\begin{theorem}[intersection of closed sets]
\begin{align*}
\left(
Q \subseteq \cc{C}{C\mathrm{\ closed}}
\middle)
\implies
\middle(
\middle(
\intersection Q
\middle)
\mathrm{\ closed}
\vphantom\union
\right)
\end{align*}
\Proof
The complement of $\intersection Q$ in the fixed ambient space is
the union of the complements of the members of $Q$. Each complement
is open, so their union is open and the intersection is closed.
\DoneNoPar
\end{theorem}

\begin{theorem}[intersection of convex sets]
\begin{align*}
\left(
Q \subseteq \cc{C}{C\mathrm{\ convex}}
\middle)
\implies
\middle(
\middle(
\intersection Q
\middle)
\mathrm{\ convex}
\vphantom\union
\right)
\end{align*}
\Proof
If $u,v\in\intersection Q$ and $0\leqslant t\leqslant1$, then
$tu+(1-t)v$ belongs to every convex member of $Q$. It therefore belongs
to the intersection, which is convex.
\DoneNoPar
\end{theorem}

\begin{theorem}[intersection of balanced sets]
\begin{align*}
\left(
Q \subseteq \cc{B}{B\mathrm{\ balanced}}
\middle)
\implies
\middle(
\middle(
\intersection Q
\middle)
\mathrm{\ balanced}
\vphantom\union
\right)
\end{align*}
\Proof
If $v\in\intersection Q$ and a scalar $a$ satisfies $|a|\leqslant1$,
then $av$ belongs to every balanced member of $Q$. It therefore belongs
to the intersection, which is balanced.
\DoneNoPar
\end{theorem}

\begin{theorem}[intersection of Dynkin systems/$\lambda$--systems]
\begin{align*}
\left(
\middle(
Q \subseteq \cc{\lambda}{\lambda\mathrm{\ Dynkin}}
\middle)
\land
\middle(
\middle(
\union\intersection Q
\middle)
\in \intersection Q
\middle)
\middle)
\implies
\middle(
\middle(
\intersection Q
\middle)
\mathrm{\ Dynkin}
\vphantom\union
\right)
\end{align*}
\Proof
Use the formulation of a Dynkin system by containment of its carrier,
closure under differences of nested members, and closure under countable
disjoint unions. The two closure properties pass to the intersection.
The displayed additional hypothesis supplies its carrier, completing
these three conditions.
\DoneNoPar
\end{theorem}

\begin{example}[Why a common measurable carrier matters]
If every sigma-algebra in $Q$ has the same carrier $X$, then
$X\in\sigma$ for every $\sigma\in Q$. Hence
$X\in\mathcal I=\bigcap Q\subseteq\power X$, which gives
$\bigcup\mathcal I=X\in\mathcal I$. Thus the additional hypothesis
in the sigma-algebra intersection theorem is automatic on a common
carrier. The same argument applies to Dynkin systems.

It is not automatic for varying carriers. For this example use the
classical real line and countable choice, so a countable union of
countable sets is countable. Let $Y=[0,\infty)$ and put
\[
\begin{split}
\Sigma&=\{A\subseteq\mathbb R:
      A\text{ is countable or }\mathbb R\setminus A
                         \text{ is countable}\},\\
\Lambda&=\power Y.
\end{split}
\]
The family $\Sigma$ is a sigma-algebra on $\mathbb R$:
complements interchange the two alternatives, and a countable union
is countable if all members are countable; if one member has countable
complement, the union has countable complement as well. The family
$\Lambda$ is a sigma-algebra on $Y$.
Since $\mathbb R\setminus Y$ is uncountable, no subset of $Y$ has
countable complement in $\mathbb R$. Consequently
\[
\Sigma\cap\Lambda=\{A\subseteq Y:A\text{ is countable}\}.
\]
Every singleton of $Y$ belongs to this intersection, so its union is
$Y$. But $Y$ is uncountable and is not a member. The intersection
therefore fails the carrier axiom, although it is closed under
countable unions and relative differences. Since sigma-algebras are
Dynkin systems, the same pair also shows why the carrier hypothesis
in that intersection theorem cannot be dropped. This is a difference
between fixed-carrier and recovered-carrier settings, not a failure of
the usual intersection theorem for sigma-algebras on one fixed space.
\end{example}

\subsection{Topological properties}
This subsection proposes a local classification by ``hereditary'' and
``weakly hereditary'' properties. The intended inheritance relation is
still to be defined: it must specify which objects are compared and which
operation or passage between them preserves a property. The preceding
intersection theorems provide one possible setting for that question.

The proposed classification treats being a topology as hereditary, being
a closed set as weakly hereditary, and being an open set as failing even
weak heredity. Compactness, separability, and connectedness are also
listed as candidate hereditary properties. These examples record the
intended direction of the investigation; their precise statements and
proofs depend on the missing definition of inheritance.

For terminological context, conventional topology distinguishes
invariance under homeomorphism from inheritance by subspaces, and may
further restrict attention to closed subspaces. The local terms here
need an explicit definition before they can be compared with those
conventions. In particular, the proposed examples should be read within
that eventual local definition.

\paragraph{A precise preservation result for continuous images.}
The following supplies one definite preservation statement for three of
the named properties, without identifying it with the proposed hierarchy.
A space is \emph{separable} if it has a countable dense subset, where
density means meeting every nonempty open set. A space is
\emph{connected} if it cannot be partitioned into two disjoint nonempty
open sets. Compactness has its open-cover meaning, stated below.

\begin{proposition}[Continuous-image preservation]
\label{prop:continuous_image_properties}
Let $f:X\to Y$ be continuous, and give $f[X]$ its subspace topology.
If $X$ is compact, separable, or connected, respectively, then $f[X]$
has that same property. No separation axiom or Choice assumption is
needed.
\end{proposition}
\Proof
Compactness is the continuous-image assertion proved in
Theorem~\ref{thm:compactness_transport}. If $D\subseteq X$ is countable
and dense, then $f[D]$ is countable. Every nonempty relatively open
$U\subseteq f[X]$ has nonempty open inverse image in $X$. That inverse
image meets $D$, so $U$ meets $f[D]$. Thus $f[D]$ is dense in $f[X]$.
Countability of the image uses no selection of preimages: an enumeration
of a nonempty countable $D$, composed with $f$, enumerates its image;
the empty case is immediate.

If $f[X]$ were partitioned into two nonempty disjoint relatively open
sets, their inverse images would be nonempty disjoint open sets
partitioning $X$. This contradicts connectedness of $X$. Empty spaces
cause no exception to any of these statements.
\Done

The operation matters. Connectedness does not pass even to closed
subspaces: $\{0,1\}$ is a disconnected closed subspace of the connected
interval $[0,1]$. For completeness, a separation of an interval would
give points $a<b$ in opposite open parts. The supremum $c$ of the first
part within $[a,b]$ can belong to neither part: if it belongs to the
first, openness and $c<b$ give a larger point in that part; if it belongs
to the second, openness and $c>a$ exclude first-part points immediately
below $c$, contradicting its supremum property. This proves the asserted
connectedness. The two endpoints themselves have disjoint open singleton
neighbourhoods in their subspace. Thus continuous-image preservation does
not determine a classification by subspaces or intersections. Defining
the proposed hierarchy still requires specifying the permitted operation
and whether its carrier is fixed or recovered.

\iftrue
\pagebreak
\subsection{Topology in older notation}
This subsection records the same neighbourhood--interior relationship
using $\nghbset_\topology(x)$ for the neighbourhoods of $x$. The symbol
$\mbox{Segm}_\sigma$ denotes arbitrary unions of the displayed interval
family; here the subscript is a union-closure notation, not a
sigma-algebra. The relation and operator interpretations must be kept
distinct in the inverse notation below.
\begin{definition}[Basic intervals of the extended real line]
\[
\mbox{Segm} = \ccc{\beta}{[-\infty;\beta)}{\beta\in\real}\cup \ccc{\alpha}{(\alpha;+\infty]}{\alpha\in\real}\cup \ccc{\alpha,\beta}{(\alpha;\beta)}{\alpha\in\real,\beta\in\real}
\]
\end{definition}
\begin{definition}[Topology of the extended real line]
\[
\topology_{[-\infty;+\infty]} = \mbox{Segm}_\sigma
\]
\end{definition}
\begin{definition}[neighbourhoods]
\[
\nghbset_\topology(x) = \cc{S \subseteq \union \topology}{ \exists  U\in \topology,\  x\in U \subseteq S }
\]
\end{definition}
\begin{definition}[interior]
Interior of a set V in a topology $\topology$.
\[
\interior_\topology(V) = \union
\separation{W}{\topology}{W\subseteq V}
\]
\end{definition}
\begin{theorem}[Converse point--set relations]
Put $E=\union\topology$ and restrict the set arguments to $\power E$.
Define the incidence relations
\[
\begin{aligned}
I&=\{(x,V)\in E\times\power E:x\in\interior(V)\},\\
N&=\{(V,x)\in\power E\times E:V\in\nghbset(x)\}.
\end{aligned}
\]
Then $I=N^{-1}$ and $N=I^{-1}$. This is the precise relational
meaning of the shorthand that interior and neighbourhood incidence
are converse; it does not assert inverse functions.
\end{theorem}
\Proof
For $x\in E$ and $V\subseteq E$, membership in the union defining
$\interior(V)$ means that some open $U$ satisfies $x\in U\subseteq V$.
This is exactly $V\in\nghbset(x)$. Thus $(x,V)\in I$ if and only if
$(V,x)\in N$, proving both identities.
The restriction is essential: for $E=\{0\}$ with topology
$\{\emptyclass,E\}$, the set $V=\{0,1\}$ has $0\in\interior(V)$
but is not an allowed neighbourhood, since $V\nsubseteq E$.
\Done
\begin{theorem}[Common neighbourhoods and interior]
Here $A,V\subseteq E=\union\topology$, and
$\nghbset_\topology\langle A\rangle$ means
$\{S\subseteq E\mid\forall x\in A,\ S\in\nghbset_\topology(x)\}$.
For $A=\emptyclass$ this is $\power E$, so the statement also covers
that case without using an unrestricted empty intersection.
\[
V\in \nghbset_\topology\langle A\rangle \iff A \subseteq \interior(V)
\]
\end{theorem}
\Proof
The left side says that $V$ is a neighbourhood of every $x\in A$.
By the preceding equivalence, this says precisely that every such $x$
lies in $\interior(V)$. Both conditions are vacuous for $A=\emptyclass$.
\Done
\begin{note}[Order-theoretic interpretation]
The shorthand $\cl\dashv\mbox{id}\dashv\interior$ records two
adjunctions, with different inclusion maps understood. Closure is left
adjoint to the inclusion of closed sets into all subsets; interior is
right adjoint to the inclusion of open sets. Concretely, for closed $C$
and open $U$,
\[
\cl(S)\subseteq C\iff S\subseteq C,\qquad
U\subseteq S\iff U\subseteq\interior(S).
\]
\end{note}

\newcommand*{\blipo}{\separation{W}{\topology}{W\subseteq V}}
\renewcommand*{\blipo}{\ccc{W}{W\in\topology}{W\subseteq V}}
\renewcommand*{\blipo}{\cc{W\in\topology}{W\subseteq V}}
% \cc{W\in\topology}{W\subseteq V}
%\noindent Lemma:
\begin{theorem}[Neighbourhood--interior correspondence]
Let $E=\union\topology$, $x\in E$, and $V\subseteq E$.
The carrier restriction on $V$ is needed because neighbourhoods, as
defined here, are subsets of $E$.
\[
V\in\nghbset_\topology(x) \iff x\in \interior_\topology(V)
\]
\Proof
$(\Rightarrow)$
\[
V\in \nghbset_\topology(x) \IMPLIES \exists U\in \topology,\  x\in U \subseteq V
\tag{nghb--elim, union--elim}
\]%\hfill(\text
\[
U\in\topology\LAND U\subseteq V \IMPLIES U \in \blipo
\tag{specification--intro}
\]
\[
x\in U \LAND U \in \blipo \IMPLIES x\in \union \blipo
\tag{union--intro}
\]

\[
x\in \union \blipo \IMPLIES x\in \interior_\topology(V)
\tag{interior--intro}
\]
$(\Leftarrow)$
\[
x\in \interior_\topology(V) \IMPLIES \exists W\in \cc{U\in\topology}{U\subseteq V},\ x\in W
\tag{interior--elim, union--elim}
\]
\[
W\in \cc{U\in\topology}{U\subseteq V} \IMPLIES W\in\topology\LAND W\subseteq V
\tag{specification--elim}
\]
\[
x\in W
\LAND
W\in\topology
\LAND
W\subseteq V
\IMPLIES
\exists U\in \topology,\  x\in U \subseteq V
\tag{and--intro, exists--intro}
\]
\[
(\exists U\in \topology,\  x\in U \subseteq V)
\IMPLIES
V\in \nghbset_\topology(x)
\tag{exists--elim, nghb--intro}
\]
\DoneNoPar
\end{theorem}
\pagebreak
\subsection{Generating families and bases}
The next two definitions retain the draft's weaker covering convention:
every open set must be expressible using the family, but the family is
not yet required to lie in $\mathcal{T}$. To obtain a base or subbase
consisting of open sets, add $\mathfrak{B}\subseteq\mathcal{T}$ or
$\mathfrak{G}\subseteq\mathcal{T}$, respectively. For finite
intersections in the generator definition, the empty intersection means
the fixed carrier $E=\union\mathcal{T}$.
%\iffalse ???
\begin{definition}[base]
A set $\mathfrak{B}$ is a base of $\mathcal{T}$ iff\\
%any open set o
$\forall U\in\mathcal{T},\ \exists W\in\power\mathfrak{B},\ U=\union W$\\
or\\
$\mathcal{T}\subseteq \cc{\union W }{  W\in\power\mathfrak{B}}$\\
or\\
$\mathcal{T}$ is a subset of unions of subsets of $\mathfrak{B}$.
\end{definition}
%\fi

%\iffalse ???
\begin{definition}[generator]
A set $\mathfrak{G}$ is a generator of $\mathcal{T}$ iff
\[
\left(\forall U\in\mathcal{T},\ \exists W\in\power\cc{\intersection Q }{ (Q\subseteq \mathfrak{G}) \land (Q \mathrm{\ finite}) },\ U=\union W
\right)
\]
or
\[
\mathcal{T}
\subseteq
\cc{\union W }{  W\in\power\cc{\intersection Q }{ (Q\subseteq \mathfrak{G}) \land (Q \mathrm{\ finite}) }}
\]
or $\mathcal{T}$ is a subset of unions of subsets of intersections of finite subsets of $\mathfrak{G}$.
\end{definition}
With the open-set condition just stated, a generator is also called a
subbase. This is the second use of that word in this chapter; unlike the
earlier intersection-closed family, it need not be closed under intersections.
%\fi

\begin{definition}[The refinement preorder on families]
For any two families $\mathcal{C} \text{ and } \mathcal{F}$ declare that $\mathcal{C} \leq \mathcal{F}$ if and only if for every $C \in \mathcal{C}$ there exists some $F \in \mathcal{F} \text{ such that } F \subseteq C,$ in which case it is said that $\mathcal{C}$ is coarser than $\mathcal{F}$ and that $\mathcal{F}$ is finer than (or subordinate to) $\mathcal{C}$.
\[
(\mathcal{C} \leq \mathcal{F})
\IFF
(\forall C \in \mathcal{C},\ \exists F \in \mathcal{F},\  F \subseteq C)
\]
\end{definition}

\begin{notation}[generator/base/topology at $x$]
\[
\mathfrak{B}(x)\ceq \cc{U\in\mathfrak{B}}{x\in U}
\]
\end{notation}

\begin{definition}[Pointwise refinement of bases]
Consistently with the preceding direction of $\leq$, base
$\mathfrak{B}_1$ is coarser than base $\mathfrak{B}_2$ iff
\[
\forall x ,\  (\mathfrak{B}_1(x) \leq \mathfrak{B}_2(x))
\]
\end{definition}

%\pagebreak
\subsection{Alexandrov's correspondence between topology and order}
Topologies closed under arbitrary intersections correspond to preorders.
The result connects the point--neighbourhood presentation with the order
theory developed earlier; see \cite[Section 1.6]{MayPishevarFiniteSpaces}.
We use the convention that open sets are \emph{upper} sets. Reversing
the preorder gives the lower-set convention used in that reference.

Fix a set $X$; no axiom of Choice or topological separation condition is used.
All intersections
below are relative to $X$, including
the empty intersection $X$; no unrestricted class intersection is intended.
A topology $\mathcal T$ on $X$ is \emph{Alexandrov} if every intersection
of its open sets is open. In that case each point has a smallest open
neighbourhood
\[
U_x=\bigcap\{U\in\mathcal T:x\in U\}.
\]
The family being intersected contains $X$, and $x\in U_x$.

\begin{theorem}[Alexandrov correspondence]
\label{thm:alexandrov_correspondence}
There is a bijection between Alexandrov topologies on $X$ and preorders
on $X$. Its two directions are
\[
x\preceq_{\mathcal T}y
\quad\Longleftrightarrow\quad
\forall U\in\mathcal T,\ x\in U\implies y\in U,
\]
and
\[
\mathcal T_{\preceq}
=\{U\subseteq X:\forall x\in U\ \forall y\in X,
                         \ x\preceq y\implies y\in U\}.
\]
In the corresponding topology, $U_x=\{y\in X:x\preceq y\}$.
\end{theorem}
\Proof
For any topology, the displayed relation is reflexive. If
$x\preceq_{\mathcal T}y$ and $y\preceq_{\mathcal T}z$, every open set
containing $x$ contains $y$ and then $z$, proving transitivity.
Every open set is upper for this relation. If the topology is Alexandrov,
then $y\in U_x$ is exactly $x\preceq_{\mathcal T}y$. Consequently an
upper subset $A$ is the union $\bigcup_{x\in A}U_x$, so it is open.
Thus the preorder recovers precisely the original open sets.

Conversely, upper subsets for a preorder include $\emptyclass$ and $X$
and are closed under all unions and all intersections. They therefore
form an Alexandrov topology. The principal upper set
$\uparrow x=\{y:x\preceq y\}$ is upper by transitivity and contains
$x$ by reflexivity. Every upper set containing $x$ contains
$\uparrow x$, so this is its smallest open neighbourhood. It follows
that the relation recovered from the topology is the original preorder.
Both constructions are inverse, proving the bijection, also when $X$
is empty.
\Done

\begin{proposition}[Continuous maps and separation]
\label{prop:alexandrov_maps}
A map between Alexandrov spaces is continuous if and only if it is
monotone for their corresponding preorders. An Alexandrov space is
$T_0$ if and only if its preorder is antisymmetric. Here $T_0$ means
that distinct points can be distinguished by an open set containing
exactly one of them.
\end{proposition}
\Proof
Let $f:X\to Y$ be continuous and $x\preceq_X x'$. For each open
$V\subseteq Y$ containing $f(x)$, its inverse image is an open set
containing $x$, hence contains $x'$. Thus $f(x)\preceq_Y f(x')$.
Conversely, if $f$ is monotone and $V$ is upper in $Y$, then
$f^{-1}[V]$ is upper in $X$, and so is open. This is continuity.

The two relations $x\preceq y$ and $y\preceq x$ hold together
exactly when $x$ and $y$ belong to the same open sets. The $T_0$
condition therefore says precisely that this can happen only when
$x=y$, which is antisymmetry.
\Done

\begin{corollary}[Finite spaces and the Hausdorff boundary]
Every finite topological space is Alexandrov. An Alexandrov Hausdorff
space is discrete; every discrete space is Alexandrov and Hausdorff.
\end{corollary}
\Proof
On a finite carrier there are only finitely many distinct open sets,
so any intersection reduces to a finite one. In a Hausdorff space,
for distinct $x,y$ there is an open neighbourhood of $x$ missing $y$.
Its associated preorder is therefore equality. The correspondence then
makes every subset open. The converse follows directly from the
definitions, including for the empty space.
\Done

For example, the preorder $0\preceq1$ on $\{0,1\}$ gives
$\mathcal T=\{\emptyclass,\{1\},\{0,1\}\}$, with
$U_0=\{0,1\}$ and $U_1=\{1\}$. This space is $T_0$ but not
Hausdorff. The indiscrete topology instead makes every pair of points
comparable in both directions. These examples show why a preorder,
rather than automatically a partial order, is the correct starting point.
The usual topology on $\mathbb R$ is not Alexandrov: the intersection
of the open intervals $(-1/n,1/n)$ over positive integers $n$ is
$\{0\}$, which is not open. Thus arbitrary intersections of open
sets remain a substantive additional hypothesis.
\pagebreak
\subsection{The product topology}
Let $(X_i,\Tau_{X_i})_{i\in I}$ be a set-indexed family of topological
spaces. The product topology on
\[
X\ceq \prod_{i\in I}X_i
\]
is the coarsest topology for which the projections
\[
\pi_i:X\to X_i
\]
are continuous for all $i\in I$.
("initial topology")
\[
\mathcal{A}\eqdef\union \cc{\cc{\pi^{-1}_i[U]}{U\in\Tau_{X_i}}}{i\in I}
\]
In fact $\mathcal{A}$ is a subbasis.
\[
\Tau_\mathcal{A}\eqdef\intersection\cc{\Tau}{(\Tau\mbox{ topology on }X)\land (\mathcal{A}\subseteq\Tau) }
\]

\begin{proposition}[Finite-coordinate basis and the product property]
\label{prop:product_topology_basis}
Let $\mathcal B$ consist of the sets
\[
\bigcap_{i\in J}\pi_i^{-1}[U_i],
\qquad J\subseteq I\text{ finite},\quad U_i\in\Tau_{X_i}.
\]
The intersection for $J=\emptyclass$ is $X$. Then the product topology
is exactly the family of unions of members of $\mathcal B$. For any
topological space $Y$ and any function $f:Y\to X$, the following are
equivalent:
\begin{enumerate}
\item $f$ is continuous;
\item every coordinate map $\pi_i\circ f:Y\to X_i$ is continuous.
\end{enumerate}
These assertions require no Choice assumption and include empty products
and products with empty carrier.
\end{proposition}
\Proof
The family $\mathcal B$ contains $X$ and is closed under binary
intersections. Indeed, combine the two finite sets of restricted
coordinates, and at a repeated coordinate replace the two open sets
by their intersection. Its unions therefore contain $\emptyclass,X$
and are closed under arbitrary unions. The intersection of two such
unions is the union of all pairwise intersections of their members,
so these unions form a topology $\mathcal S$.

Every subbasic inverse image belongs to $\mathcal B$, hence to
$\mathcal S$. Conversely, any topology containing $\mathcal A$
contains its finite intersections and all their unions. It therefore
contains $\mathcal S$. Thus $\mathcal S$ equals the displayed
intersection defining $\Tau_{\mathcal A}$; that intersection is over
a nonempty set of topologies, since the discrete topology $\power X$
is one of them. This proves the basis assertion and the coarsest-topology
characterization simultaneously. In particular each $\pi_i$ is continuous.

If $f$ is continuous, inverse images under $\pi_i\circ f$ are obtained
by taking two continuous inverse images, so each coordinate map is
continuous. Conversely, suppose every coordinate map is continuous.
The inverse image under $f$ of a member of $\mathcal B$ is the finite
intersection of the corresponding open coordinate inverse images.
It is open, including when the intersection is empty and its inverse
image is $Y$. Since inverse images commute with unions, $f$ is
continuous on every product-open set.

When $I=\emptyclass$, $X$ is the singleton containing the empty
function, $\mathcal B=\{X\}$, and every map into it is continuous.
When $X=\emptyclass$, the only topology on it is $\{\emptyclass\}$;
a function $Y\to X$ exists only for $Y=\emptyclass$, and both
conditions then hold. The construction never assumes that a product
of nonempty carriers has a point.
\Done

For a specified point $x\in X$, the basis assertion says more locally:
$N\subseteq X$ is a neighbourhood of $x$ exactly when it contains
some finite-coordinate basic set containing $x$. In a binary product
$E\times F$, this is equivalently the condition
\[
\exists A\in\mathrm{nghb}_E(a),\quad
\exists B\in\mathrm{nghb}_F(b),\quad A\times B\subseteq N,
\qquad x=(a,b).
\]
Indeed an open basic rectangle gives such $A,B$. Conversely, choose
open neighbourhoods $U,V$ with $a\in U\subseteq A$ and
$b\in V\subseteq B$; then $U\times V$ is an open neighbourhood
inside $N$. Thus the earlier function $\mathrm{pfb}((a,b))$ really
is a filter base generating precisely the product neighbourhood filter.
It is nonempty, all its members contain $(a,b)$, and intersecting two
members gives another neighbourhood rectangle. This connects the
pointwise filter construction with the open-set definition.

\fi %iftrue
%Classification based on
\subsection{Cardinality--based topological definitions}
Let $(X,\Tau)$ be a topological space and $S\subseteq X$. In this
subsection, $\mathcal{N}(x)$ denotes the neighbourhood family of $x$.
The formulas define the author's cardinality-based terminology directly.
In particular, ``limits'' requires at least two points of $S$ in each
neighbourhood; no separation assumption is implicit. The names
``isolates'' and ``border'' refer to the differences displayed below.

\[
\text{closure}(S)\eqdef\cc{x\in X}{\forall N\in\mathcal{N}(x),\card(S\cap N) \geqslant 1}
\]

\[
\text{exterior}(S) \eqdef X\setminus \text{closure}(S)
\]

\[
\text{interior}(S) \eqdef \text{exterior}(X\setminus S)
\]
\iffalse
\[
\text{interior}(S) \eqdef X\setminus \text{closure}(X\setminus S)
\]
\fi
\[
\text{limits}(S)\eqdef \cc{x\in X}{\forall N\in\mathcal{N}(x),\card(S\cap N) \geqslant 2}
\]

\[
\text{isolates}(S) \eqdef \text{closure}(S)\setminus\text{limits}(S)
\]

\[
\text{border}(S) \eqdef \text{closure}(S)\setminus\text{interior}(S)
\]

\begin{proposition}[Comparison with the ordinary derived set]
\label{prop:cardinality_limits_T1}
Define the ordinary derived set by
\[
S'=\{x\in X:\text{every neighbourhood of }x
                 \text{ meets }S\setminus\{x\}\}.
\]
In every topological space, $\text{limits}(S)\subseteq S'$.
If $X$ is $T_1$, meaning that every singleton is closed, then
$\text{limits}(S)=S'$. In that case $\text{isolates}(S)$ consists
exactly of the isolated points of the subspace $S$.
\end{proposition}
\Proof
If a neighbourhood contains two distinct points of $S$, at least one
differs from $x$. This proves the first inclusion without separation
assumptions. Conversely, suppose $X$ is $T_1$, $x\in S'$, and $N$
is a neighbourhood of $x$. If $x\in S$, then $N$ contains $x$ and
a different point of $S$, as required. If $x\notin S$, choose
$y\in N\cap S$. The set $N\setminus\{y\}$ is still a
neighbourhood of $x$: it contains the intersection of an open
neighbourhood inside $N$ with the open set $X\setminus\{y\}$.
Since $x\in S'$, this smaller neighbourhood meets $S$, giving a
second point different from $y$. Thus every $N$ contains at least
two points of $S$.

If $x\in\text{closure}(S)\setminus S$, every neighbourhood meets
$S\setminus\{x\}=S$, so $x\in S'$. Hence
$\text{closure}(S)\setminus S'=S\setminus S'$.
For $x\in S$, failure to belong to $S'$ means that some neighbourhood
meets $S$ only in $x$, equivalently $\{x\}$ is open in the subspace
$S$. This proves the assertion about isolated points.
\Done

The separation assumption cannot simply be omitted. In the indiscrete
space $X=\{a,b\}$ with $a\ne b$, take $S=\{a\}$. The only
neighbourhood of either point is $X$, so $S'=\{b\}$ whereas
$\text{limits}(S)=\emptyclass$. Here the author's
$\text{isolates}(S)=X$ even contains a point outside $S$.
The two conventions therefore describe different notions in this setting.

\iffalse
\begin{tabular}{|l|l|}
\hline
Set & \text{Column A} & \text{Column B} \\
Cx & y \\
\end{tabular}
\fi
%\todo
\pagebreak



\subsection{Compactness and changes of topology}
In this subsection all spaces are set sized. A space is \emph{compact}
if every cover by open sets has a finite subcover. A subset is compact
when it is compact with its subspace topology. These definitions do
not impose a separation axiom.

\begin{theorem}[Transport of compactness]
\label{thm:compactness_transport}
A closed subset of a compact space is compact, and the continuous
image of a compact space is compact. Every compact subset of a
Hausdorff space is closed. Consequently a continuous bijection from
a compact space to a Hausdorff space is a homeomorphism.
\end{theorem}
\Proof
Let $F$ be closed in a compact space $X$. Write each member of a
relative open cover of $F$ as $F\cap U$ for an ambient open set $U$.
One need not choose representatives simultaneously: take all open
$U$ whose intersection with $F$ is a member of the given cover.
These sets together with $X\setminus F$ cover $X$. A finite subcover
therefore yields a finite subcover of $F$ from the original cover.

For a continuous $f:X\to Y$, the inverse images of an open cover of
$f[X]$ cover $X$ by open sets. A finite subcover pulls forward to a
finite subcover of $f[X]$.

Let $K\subseteq Y$ be compact, with $Y$ Hausdorff, and fix
$y\in Y\setminus K$. Consider all open sets $U$ meeting $K$ for
which there is an open neighbourhood $V$ of $y$ disjoint from $U$.
Hausdorff separation makes these $U$ cover $K$. Finitely many
$U_1,\ldots,U_m$ cover $K$; choose their corresponding $V_i$.
The finite intersection of the $V_i$ is a neighbourhood of $y$
disjoint from $K$. For $K=\emptyclass$ the conclusion is immediate.
Thus $Y\setminus K$ is open, proving closedness. Only finitely many
choices of witnesses were used.

Finally, a continuous bijection $f:X\to Y$ as stated sends every
closed subset of $X$ to a compact, hence closed, subset of $Y$.
This is exactly continuity of its inverse, since the inverse image
under $f^{-1}$ of a closed subset $F\subseteq X$ is $f[F]$.
\Done

The Hausdorff hypothesis is substantive. On a two-point set, the
identity from the discrete topology to the indiscrete topology is a
continuous bijection between compact spaces. Its inverse is not
continuous. Thus compactness alone does not recover the original
topology from a continuous bijection.

\subsection{Compactness in metric spaces}
This subsection uses ordinary set-sized metric spaces and classical
countable choice. A space is \emph{compact} when every open cover has
a finite subcover; it is \emph{totally bounded} when, for every
$\varepsilon>0$, finitely many open balls of radius $\varepsilon$
cover it. These are distinct definitions.

\begin{theorem}[The metric compactness criterion]
For a metric space $X$, the following conditions are equivalent:
\begin{enumerate}
\item $X$ is compact;
\item every sequence in $X$ has a convergent subsequence with limit in $X$;
\item $X$ is complete and totally bounded.
\end{enumerate}
\end{theorem}
\Proof
The empty space satisfies all three conditions, so assume $X$ nonempty.
Compactness gives total boundedness by applying the cover property to
all balls of a fixed radius. If $(x_n)$ is Cauchy, the closed sets
$K_n=\overline{\{x_j:j\geq n\}}$ are nonempty and nested.
They have nonempty intersection: otherwise their open complements
would cover $X$, and a finite subcover would make one $K_n$ empty.
Choose $x$ in the intersection. Given $\varepsilon>0$, a sufficiently
late tail lies in the closed ball of radius $\varepsilon/2$ about one
of its terms; its closure contains $x$. The triangle inequality then
places all sufficiently late $x_n$ within $\varepsilon$ of $x$.
Thus compactness implies completeness and total boundedness.

Assume completeness and total boundedness. By countable choice, fix in
advance, for every $k\geq1$, a finite ordered list of radius $2^{-k}$
balls covering $X$. Start with all sequence indices. At stage $k$,
retain those indices in the first ball of the $k$th list containing
infinitely many of the retained terms; such a ball exists because the
list is finite. Choose $n_k$ as the least retained index greater than
$n_{k-1}$, starting with $n_0=0$. Once the lists are fixed these are
uniquely specified recursive operations, not dependent choices.
For $j,l\geq k$,
the chosen terms lie in a common ball of radius $2^{-k}$, so their
distance is at most $2^{1-k}$. The subsequence is Cauchy and hence
converges in $X$.

Finally assume the subsequence condition. If total boundedness failed
at some radius $r>0$, finite induction would produce, for every positive
integer $n$, an ordered $n$-tuple with pairwise distances at least $r$:
at each finite stage some point lies outside the finitely many radius
$r$ balls about the points already obtained. Countable choice selects
one such tuple for each $n$. Concatenate these finite tuples to obtain
a sequence enumerating their union $S$. No finite collection of radius
$r/2$ balls covers $S$, since each such ball contains at most one point
of any selected tuple, and tuples have arbitrarily large lengths.
Starting with the first enumerated point, recursively choose the first
point in the enumeration outside the radius $r/2$ balls about the
previously chosen points. This rule is uniquely specified and always
possible by the preceding observation. It yields a sequence with pairwise
distances at least $r/2$, and hence with no convergent subsequence.
This contradiction proves total boundedness using only countable choice.
For an open cover $\mathcal U$, there is a number $\delta>0$ such
that every ball $B(x,\delta)$ is contained in some member of
$\mathcal U$. If not, choose $x_n$ whose ball of radius $1/n$ is
contained in no member. A convergent subsequence has limit $x$ in
some open $U\in\mathcal U$. Choose $r>0$ with $B(x,2r)\subseteq U$.
Eventually both $d(x_{n_k},x)<r$ and $1/n_k<r$, contradicting the
choice of $x_{n_k}$. Now finitely many radius $\delta$ balls cover
$X$ by total boundedness, and containing cover members give a finite
subcover. Hence $X$ is compact.
\Done

\begin{corollary}[Relative norm compactness in a Banach space]
A subset of a Banach space has compact norm closure if and only if it
is totally bounded in the norm metric.
\end{corollary}
\Proof
A compact closure is totally bounded, and a subset inherits that
property (replace a covering centre outside the subset by a point of
an intersecting smaller ball). Conversely, the closure of a totally
bounded set is totally bounded: a finite cover by radius
$\varepsilon/3$ balls covers its closure by radius $\varepsilon$ balls.
The closure is complete, being closed in a complete space. Apply the
metric criterion.
\Done

For example, the closed unit ball of the infinite-dimensional Hilbert
space $\ell^2$ is complete and bounded, but not compact: its coordinate
vectors have pairwise distance $\sqrt2$, so no finite collection of
balls of radius less than $\sqrt2/2$ covers them. Total boundedness,
not boundedness alone, is the needed condition. The equivalence with
sequential compactness proved here is a metric theorem; it must not be
applied to an arbitrary weak topology without a separate argument.

\subsection{Contractions on complete metric spaces}
The next result uses completeness directly and does not assume
compactness or total boundedness. A contraction here has one fixed
constant $q<1$ controlling all pairs of points.

\begin{theorem}[Banach's contraction principle and error bounds]
\label{thm:banach_contraction}
Let $(X,d)$ be a nonempty complete metric space and let $T:X\to X$
satisfy
\[
d(Tx,Ty)\leq q\,d(x,y)\qquad(x,y\in X),\qquad 0\leq q<1.
\]
There is a unique $p\in X$ with $Tp=p$. For every initial point
$x_0\in X$, the iteration $x_{n+1}=Tx_n$ converges to $p$.
For every $n\geq1$ it satisfies the a priori and a posteriori bounds
\[
d(x_n,p)\leq\frac{q^n}{1-q}d(x_1,x_0),\qquad
d(x_n,p)\leq\frac{q}{1-q}d(x_n,x_{n-1}).
\]
More generally, every $x\in X$ satisfies the residual estimate
\[
d(x,p)\leq\frac{d(x,Tx)}{1-q}.
\]
In particular this last estimate also applies to the initial point.
\end{theorem}
\Proof
If $q=0$, all values of $T$ coincide. Write $p=Tx_0$ for any
$x_0\in X$. Then $Tp=p$, and $x_n=p$ for every $n\geq1$.
The displayed iteration bounds are both zero, as required. Any fixed
point equals this same value of $T$, so uniqueness also holds.

Now suppose $0<q<1$ and fix $x_0$. Natural-number recursion uniquely
defines the iteration; after fixing $x_0$ no arbitrary selections are
made. Put $D=d(x_1,x_0)$. Induction gives
\[
d(x_{k+1},x_k)\leq q^kD\qquad(k\geq0).
\]
For integers $m>n\geq0$, the triangle inequality and the finite
geometric sum give
\[
d(x_m,x_n)\leq\sum_{k=n}^{m-1}q^kD
=\frac{q^n(1-q^{m-n})}{1-q}D
\leq\frac{q^n}{1-q}D.
\]
Since $q^n\to0$, the sequence is Cauchy; this includes $D=0$,
when all iterates are already equal. Completeness gives $x_n\to p$
for some $p\in X$. The estimate
\[
d(Tp,p)\leq d(Tp,Tx_n)+d(x_{n+1},p)
\leq q\,d(p,x_n)+d(x_{n+1},p)
\]
tends to zero, proving $Tp=p$. If $z$ is another fixed point, then
$d(p,z)=d(Tp,Tz)\leq qd(p,z)$, hence $d(p,z)=0$ and $p=z$.
This uniqueness also shows that all initial points lead to the same
limit. Letting $m\to\infty$ in the preceding geometric bound proves
the a priori estimate.

For either $q=0$ or $0<q<1$ and any $x\in X$,
\[
d(x,p)\leq d(x,Tx)+d(Tx,Tp)
\leq d(x,Tx)+q\,d(x,p).
\]
Rearranging proves the residual estimate. For $n\geq1$, contraction gives
\[
d(x_n,Tx_n)\leq q\,d(x_{n-1},x_n).
\]
Applying the residual estimate to $x_n$ now gives the a posteriori bound.
Every denominator is positive; the
case $q=0$ uses no expression $0^0$ in the stated iteration bounds.
\Done

\begin{corollary}[Stability of the fixed point]
Under the theorem's assumptions, suppose another self-map $S:X\to X$
has a fixed point $z$ and $d(Tz,Sz)\leq\eta$ for some $\eta\geq0$.
Then
\[
d(p,z)\leq\frac{\eta}{1-q}.
\]
In particular the hypothesis holds if $d(Tx,Sx)\leq\eta$ for every
$x\in X$ and $S$ has a fixed point.
\end{corollary}
\Proof
Since $Sz=z$, the residual at $z$ is
$d(z,Tz)=d(Sz,Tz)\leq\eta$. Apply the residual estimate.
\Done

\begin{example}[Completeness and strict contraction are essential]
On the incomplete space $(0,1)$ with its usual metric, $T(x)=x/2$
is a contraction with $q=1/2$ and maps the space into itself. Its
only possible fixed point is zero, which is outside the space; its
iterates converge to that missing endpoint in $\mathbb R$.
On the complete space $\mathbb R$, the translation $T(x)=x+1$
satisfies $d(Tx,Ty)=d(x,y)$ but has no fixed point. Thus the estimate
with $q=1$ is insufficient. Nonemptiness matters too: the unique
self-map of the empty complete metric space satisfies every displayed
contraction inequality vacuously but has no fixed point.
\end{example}

\subsection{Proper ultrafilters and compact products}
This subsection concerns set-sized topological spaces and uses classical
set theory with Choice. Its filters are \emph{proper filters on a fixed
carrier}; this qualification distinguishes them from any earlier convention
allowing empty filters or varying carriers. These results will supply the
compactness arguments for weak topologies, where the metric subsequence
criterion alone is insufficient.

\begin{definition}[Proper filters and ultrafilters]
A proper filter $\mathcal F$ on a set $X$ is a family of subsets of $X$
containing $X$, excluding $\emptyclass$, closed under finite intersections,
and upward closed among subsets of $X$. An ultrafilter is a proper filter
maximal under inclusion among proper filters on the same carrier.
A family $\mathcal A\subseteq\power X$ has the finite intersection
property when every finite intersection of its members is nonempty;
the empty intersection is $X$.
\end{definition}

\begin{theorem}[Ultrafilter extension and the dichotomy]
\label{thm:ultrafilter_extension}
Every family with the finite intersection property is contained in an
ultrafilter on $X$. A proper filter $\mathcal U$ is an ultrafilter exactly
when, for every $A\subseteq X$, exactly one of $A$ and $X\setminus A$
belongs to $\mathcal U$.
\end{theorem}
\Proof
The supersets of finite intersections from $\mathcal A$ form a proper
filter $\mathcal F_0$. Order its proper-filter extensions by inclusion.
The union of a nonempty chain is a proper filter: any finitely many of
its members lie together in one chain member, and none of those filters
contains the empty set. For the empty chain, use $\mathcal F_0$ as an
upper bound. Zorn's lemma gives a maximal extension $\mathcal U$.

If $A\notin\mathcal U$, adjoining $A$ cannot generate a proper filter,
by maximality. Thus some $F\in\mathcal U$ has $F\cap A=\emptyclass$;
otherwise the finite intersections from $\mathcal U\cup\{A\}$ would
generate a proper extension. Upward closure gives
$X\setminus A\in\mathcal U$. Both alternatives cannot occur, since
their intersection is empty. Conversely, if the dichotomy holds and a
proper filter extending $\mathcal U$ contained a new set $A$, it would
also contain $X\setminus A\in\mathcal U$, a contradiction. This proves
the characterization.
\Done

\begin{definition}[Ultrafilter convergence and direct image]
An ultrafilter $\mathcal U$ on a topological space $X$ converges to $x$
if every neighbourhood of $x$ belongs to $\mathcal U$. For any map
$f:X\to Y$, define
\[
f_*\mathcal U=\{B\subseteq Y:f^{-1}[B]\in\mathcal U\}.
\]
\end{definition}

\begin{lemma}[Transport and uniqueness of ultrafilter limits]
The family $f_*\mathcal U$ is an ultrafilter, without any continuity
assumption on $f$. If $f$ is continuous and $\mathcal U$ converges to $x$,
then $f_*\mathcal U$ converges to $f(x)$. In a Hausdorff space an
ultrafilter has at most one limit.
\end{lemma}
\Proof
Inverse images preserve the carrier, the empty set, intersections,
inclusions, and relative complements. The filter axioms and the dichotomy
therefore pass to $f_*\mathcal U$. For continuity, the inverse image of
a neighbourhood of $f(x)$ is a neighbourhood of $x$, hence belongs to
$\mathcal U$. Finally, disjoint neighbourhoods of two distinct Hausdorff
points cannot both belong to a proper filter.
\Done

\begin{theorem}[Compactness characterized by ultrafilters]
\label{thm:compact_ultrafilter}
A topological space is compact if and only if every ultrafilter on its
carrier converges to some point of the space.
\end{theorem}
\Proof
Suppose $X$ is compact and $\mathcal U$ is an ultrafilter. The closed
sets $\overline F$, $F\in\mathcal U$, have the finite intersection
property: the corresponding finite intersection of the $F$ is nonempty
and contained in their closures. Compactness implies that their total
intersection is nonempty, since otherwise the open complements would
have a finite subcover. Choose $x$ in that intersection. If a neighbourhood
$N$ of $x$ did not belong to $\mathcal U$, the dichotomy would put
$X\setminus N$ in $\mathcal U$. But an open neighbourhood $O$ with
$x\in O\subseteq N$ shows that
$x\notin\overline{X\setminus N}$, a contradiction. Thus
$\mathcal U$ converges to $x$.

Conversely, if an open cover has no finite subcover, its complements
have the finite intersection property. Extend them to an ultrafilter.
A limit $x$ would lie in some open cover member $O$, putting both $O$
and $X\setminus O$ in that ultrafilter. This is impossible. Hence the
cover has a finite subcover. The empty space is covered by this argument:
it is compact and carries no proper ultrafilter.
\Done

\begin{theorem}[Compactness of set-indexed products]
\label{thm:tychonoff}
Let $(X_i)_{i\in I}$ be a set-indexed family of compact spaces. Their
product, with the product topology, is compact. Separation axioms on the
factors are not required.
\end{theorem}
\Proof
If the product $X$ is empty, it is compact. Otherwise let $\mathcal U$
be an ultrafilter on $X$. For each coordinate projection $\pi_i$, the
ultrafilter $(\pi_i)_*\mathcal U$ has a limit in $X_i$ by compactness.
Choice selects such a limit $x_i$ for each $i$; put $x=(x_i)\in X$.
Every basic neighbourhood of $x$ restricts only finitely many coordinates.
Its coordinate neighbourhoods belong to the respective image ultrafilters,
so their inverse images belong to $\mathcal U$. Their finite intersection
therefore belongs to $\mathcal U$, and upward closure handles every
neighbourhood of $x$. Thus $\mathcal U$ converges to $x$. The preceding
theorem proves compactness. For $I=\emptyclass$, the product is the
singleton consisting of the empty function, and the same finite-coordinate
argument has no conditions to impose.
\Done

In particular, a product of closed bounded real intervals is compact.
Each interval is complete as a closed subset of the complete real line
and totally bounded by finite grids, so the earlier metric criterion
proves its compactness. Product compactness here is genuinely topological;
no sequential compactness claim for an arbitrary uncountable product has
been substituted for it. The proof uses Choice explicitly, unlike the
deterministic iteration in the contraction theorem.

\ifdefined\qwertyflag
\undef\qwertyflag
\end{document}
\fi
